% =============================================================================
%  Notas de clase — Sesión 04: Representaciones Binarias
%  Programación Avanzada — Universidad Panamericana
%  Compilar: pdflatex sesion04_representaciones_binarias.tex
% =============================================================================
\documentclass[12pt, oneside]{article}
\usepackage{progavanzada}

\begin{document}

\portada{Sesión 04}{Representaciones Binarias}{2026}

\tableofcontents
\newpage

% =============================================================================
\section{Solución: el problema de Carina}
% =============================================================================

El problema de la sesión anterior: Carina escribe los enteros positivos
en orden pero su calculadora no tiene las teclas del \textbf{2} ni del
\textbf{7}. ¿Cuál es el milésimo número que escribe?

Los dígitos disponibles son $\{0, 1, 3, 4, 5, 6, 8, 9\}$ — exactamente
\textbf{ocho} símbolos. Si los renombramos como $\{0,1,2,3,4,5,6,7\}$,
Carina en realidad está contando en \textbf{base 8}.

\begin{ejemplo}[Encontrar el milésimo número]
Convertimos $1000_{10}$ a base 8 por divisiones sucesivas:
\begin{center}
\begin{tabular}{r@{\;}c@{\;}l@{\quad}l}
  1000 & $\div$ & 8 = 125 & resto \textbf{0} \\
   125 & $\div$ & 8 = 15  & resto \textbf{5} \\
    15 & $\div$ & 8 = 1   & resto \textbf{7} \\
     1 & $\div$ & 8 = 0   & resto \textbf{1} \\
\end{tabular}
\end{center}
$1000_{10} = 1750_8$. Ahora mapeamos cada dígito octal al símbolo real de Carina:
\[
  0\to0,\quad 1\to1,\quad 2\to3,\quad 3\to4,\quad
  4\to5,\quad 5\to6,\quad 6\to8,\quad 7\to9
\]
$1750_8 \;\xrightarrow{\text{remap}}\; \mathbf{1960}$
\end{ejemplo}

\begin{recuerda}
  Verificación: $1\times8^3 + 7\times8^2 + 5\times8 + 0 = 512+448+40 = 1000$. \checkmark
\end{recuerda}

% =============================================================================
\section{¿Cuántos valores caben en $n$ bits?}
% =============================================================================

Con $n$ bits existen exactamente $2^n$ patrones posibles.
Para enteros \textbf{sin signo} (solo positivos), esos patrones representan
los valores $0$ a $2^n - 1$.

\begin{center}
\begin{tabular}{ccl}
  \toprule
  \textbf{Bits} & \textbf{Patrones} & \textbf{Rango sin signo} \\
  \midrule
   8 & $2^8  = 256$                   & 0 a 255 \\
  16 & $2^{16} \approx 65{,}000$      & 0 a 65\,535 \\
  32 & $2^{32} \approx 4{,}000{,}000{,}000$ & 0 a 4\,294\,967\,295 \\
  64 & $2^{64} \approx 1.8\times10^{19}$ & 0 a $\approx$18 trillones \\
  \bottomrule
\end{tabular}
\end{center}

\begin{consejo}
  Memoriza las potencias clave: $2^8=256$, $2^{10}=1024$ (aprox.\ 1\,000),
  $2^{16}=65\,536$, $2^{32}\approx4\times10^9$. Son las que aparecen
  constantemente al declarar tipos en C++.
\end{consejo}

% =============================================================================
\section{Representación de negativos — Complemento a 2}
% =============================================================================

Para representar números negativos hay que ``partir'' el rango de $2^n$
patrones en positivos y negativos. El método que usa \textbf{toda la
computación moderna} es el \textbf{complemento a 2}.

\subsection{La idea}

Con $n$ bits dedicamos \textbf{1 bit} (el más significativo, el de más a
la izquierda) a indicar el signo:
\begin{itemize}
  \item \textbf{0} en el bit más significativo → número \textbf{positivo o cero}.
  \item \textbf{1} en el bit más significativo → número \textbf{negativo}.
\end{itemize}

Con 8 bits el reparto queda:

\begin{center}
\begin{tabular}{cc}
  \toprule
  \textbf{Patrón} & \textbf{Valor} \\
  \midrule
  \texttt{00000000} & $0$ \\
  \texttt{00000001} & $+1$ \\
  $\vdots$ & $\vdots$ \\
  \texttt{01111111} & $+127$ \\
  \texttt{10000000} & $-128$ \\
  \texttt{10000001} & $-127$ \\
  $\vdots$ & $\vdots$ \\
  \texttt{11111110} & $-2$ \\
  \texttt{11111111} & $-1$ \\
  \bottomrule
\end{tabular}
\end{center}

\begin{recuerda}
  Con $n$ bits en complemento a 2 el rango es $-2^{n-1}$ a $+2^{n-1}-1$.
  Para 8 bits: \textbf{$-128$ a $+127$}. Hay un negativo ``extra''
  porque el cero ocupa un patrón del lado positivo.
\end{recuerda}

\subsection{¿Por qué \texttt{11111111} es $-1$ y no $-127$?}

La clave: en complemento a 2, un número negativo $-x$ se almacena como
$2^n - x$.
\[
  -1 \;\text{con 8 bits}\; = 2^8 - 1 = 256 - 1 = 255 = \texttt{11111111}_2
\]
\[
  -128 \;\text{con 8 bits}\; = 2^8 - 128 = 256 - 128 = 128 = \texttt{10000000}_2
\]

\subsection{Cómo negar un número: voltear + 1}

Para obtener $-x$ a partir de $x$ (o viceversa):
\begin{enumerate}
  \item \textbf{Invierte} todos los bits (0→1, 1→0).
  \item \textbf{Suma 1} al resultado.
\end{enumerate}

\begin{ejemplo}[Calcular $-6$ en 8 bits]
\begin{center}
\begin{tabular}{ll}
  $6_{10}$              & \texttt{00000110} \\
  Invertir bits         & \texttt{11111001} \\
  Sumar 1               & \texttt{11111010} \\
\end{tabular}
\end{center}
$-6_{10} = \texttt{11111010}_2$. \\[4pt]
Verificación: $-6 = 2^8 - 6 = 250 = \texttt{11111010}_2$. \checkmark
\end{ejemplo}

\begin{advertencia}
  El rango no es simétrico: con 8 bits existe $-128$ pero \textbf{no}
  $+128$. Si intentas negar $-128$ aplicando el algoritmo, el resultado
  desborda y obtienes $-128$ de nuevo. Este es un comportamiento real
  de las CPUs y fuente de bugs históricos.
\end{advertencia}

% =============================================================================
\section{ASCII — El texto también son números}
% =============================================================================

Un archivo de texto no es más que una secuencia de bytes.
El estándar \textbf{ASCII} (\textit{American Standard Code for Information
Interchange}, 1963) asigna un número a cada carácter imprimible y de control.

\subsection{Tabla de letras mayúsculas}

\begin{center}
\begin{tabular}{cc|cc|cc}
  \toprule
  \textbf{Letra} & \textbf{Hex} &
  \textbf{Letra} & \textbf{Hex} &
  \textbf{Letra} & \textbf{Hex} \\
  \midrule
  A & 41 & J & 4A & S & 53 \\
  B & 42 & K & 4B & T & 54 \\
  C & 43 & L & 4C & U & 55 \\
  D & 44 & M & 4D & V & 56 \\
  E & 45 & N & 4E & W & 57 \\
  F & 46 & O & 4F & X & 58 \\
  G & 47 & P & 50 & Y & 59 \\
  H & 48 & Q & 51 & Z & 5A \\
  I & 49 & R & 52 &   &    \\
  \bottomrule
\end{tabular}
\end{center}

\begin{ejemplo}[``MAR'' en ASCII]
\begin{center}
\begin{tabular}{ccc}
  \toprule
  \textbf{Letra} & \textbf{Hex} & \textbf{Binario} \\
  \midrule
  M & 4D & \texttt{0100 1101} \\
  A & 41 & \texttt{0100 0001} \\
  R & 52 & \texttt{0101 0010} \\
  \bottomrule
\end{tabular}
\end{center}
El texto ``MAR'' ocupa exactamente \textbf{3 bytes} en memoria.
\end{ejemplo}

\subsection{Mayúsculas vs.\ minúsculas}

Observa el patrón binario de las mayúsculas y sus minúsculas:

\begin{center}
\begin{tabular}{ccc}
  \toprule
  & \textbf{Hex} & \textbf{Binario} \\
  \midrule
  \texttt{A} & 41 & \texttt{010\textbf{0} 0001} \\
  \texttt{a} & 61 & \texttt{011\textbf{0} 0001} \\
  \midrule
  \texttt{M} & 4D & \texttt{010\textbf{0} 1101} \\
  \texttt{m} & 6D & \texttt{011\textbf{0} 1101} \\
  \bottomrule
\end{tabular}
\end{center}

La única diferencia es el \textbf{bit 5} (contando desde 0 por la derecha):
\begin{itemize}
  \item Mayúsculas: bit 5 = \textbf{0} → patrones \texttt{010xxxxx} (0x4\_ y 0x5\_)
  \item Minúsculas: bit 5 = \textbf{1} → patrones \texttt{011xxxxx} (0x6\_ y 0x7\_)
\end{itemize}

Convertir mayúscula $\leftrightarrow$ minúscula es simplemente \textbf{activar o desactivar el bit 5},
equivalente a sumar o restar 32.

\begin{consejo}
  \texttt{'A'} + 32 = \texttt{'a'}. \quad \texttt{'Z'} + 32 = \texttt{'z'}. \\
  Esto explica por qué en C++ puedes convertir entre casos con una suma:
  \texttt{c + 32} o, más elegante, \texttt{c \textasciicircum{} 0x20}
  (XOR con el bit 5).
\end{consejo}

\subsection{Dígitos y signos de puntuación}

Los caracteres \texttt{'0'}–\texttt{'9'} tienen un patrón muy limpio:
el nibble alto siempre es \texttt{0011} (hex \texttt{3}) y el nibble bajo
es \textbf{exactamente el valor binario del dígito}.

\begin{center}
\begin{tabular}{ccc}
  \toprule
  \textbf{Carácter} & \textbf{Hex} & \textbf{Binario} \\
  \midrule
  \texttt{'0'} & 30 & \texttt{0011 \textbf{0000}} \\
  \texttt{'1'} & 31 & \texttt{0011 \textbf{0001}} \\
  \texttt{'2'} & 32 & \texttt{0011 \textbf{0010}} \\
  \texttt{'3'} & 33 & \texttt{0011 \textbf{0011}} \\
  \texttt{'4'} & 34 & \texttt{0011 \textbf{0100}} \\
  \texttt{'5'} & 35 & \texttt{0011 \textbf{0101}} \\
  \texttt{'6'} & 36 & \texttt{0011 \textbf{0110}} \\
  \texttt{'7'} & 37 & \texttt{0011 \textbf{0111}} \\
  \texttt{'8'} & 38 & \texttt{0011 \textbf{1000}} \\
  \texttt{'9'} & 39 & \texttt{0011 \textbf{1001}} \\
  \bottomrule
\end{tabular}
\end{center}

\begin{consejo}
  Para convertir \texttt{'7'} al número entero 7 basta con hacer
  \texttt{'7' \& 0x0F} (AND con los 4 bits bajos) o simplemente
  \texttt{'7' - '0'}. Ambas operaciones eliminan el prefijo \texttt{0011}.
\end{consejo}

Los signos de control y puntuación ocupan los valores más bajos y
medios de la tabla:

\begin{center}
\begin{tabular}{lcc}
  \toprule
  \textbf{Carácter} & \textbf{Hex} & \textbf{Binario} \\
  \midrule
  \texttt{NUL} (nulo \texttt{\textbackslash 0})        & 00 & \texttt{0000 0000} \\
  \texttt{TAB} (tabulador \texttt{\textbackslash t})   & 09 & \texttt{0000 1001} \\
  \texttt{LF}  (salto de línea \texttt{\textbackslash n}) & 0A & \texttt{0000 1010} \\
  \texttt{CR}  (retorno de carro \texttt{\textbackslash r}) & 0D & \texttt{0000 1101} \\
  \texttt{ESC} (escape)                                & 1B & \texttt{0001 1011} \\
  \midrule
  \texttt{' '} (espacio)   & 20 & \texttt{0010 0000} \\
  \texttt{'!'}             & 21 & \texttt{0010 0001} \\
  \texttt{'"'}             & 22 & \texttt{0010 0010} \\
  \texttt{'+'}             & 2B & \texttt{0010 1011} \\
  \texttt{'-'}             & 2D & \texttt{0010 1101} \\
  \texttt{'/'}             & 2F & \texttt{0010 1111} \\
  \midrule
  \texttt{':'}             & 3A & \texttt{0011 1010} \\
  \texttt{';'}             & 3B & \texttt{0011 1011} \\
  \texttt{'='}             & 3D & \texttt{0011 1101} \\
  \midrule
  \texttt{'\{'}            & 7B & \texttt{0111 1011} \\
  \texttt{'\}'}            & 7D & \texttt{0111 1101} \\
  \texttt{DEL} (borrar)    & 7F & \texttt{0111 1111} \\
  \bottomrule
\end{tabular}
\end{center}

\begin{recuerda}
  El espacio (\texttt{0x20} = \texttt{0010 0000}) y el \texttt{NUL}
  (\texttt{0x00}) son caracteres reales con su propio byte.
  Un archivo de texto que ``parece vacío'' puede estar lleno de espacios,
  tabuladores o saltos de línea — todos con representación binaria propia.
\end{recuerda}

\begin{tecnico}[ASCII extendido y Unicode]
  El ASCII original usa 7 bits (128 caracteres). Para incluir acentos,
  chino, árabe y el resto de los sistemas de escritura del mundo existe
  \textbf{Unicode}, que hoy define más de 149\,000 caracteres.
  La codificación más común en la web y en sistemas modernos es
  \textbf{UTF-8}, compatible con ASCII para los primeros 128 valores.\\
  \textit{Referencia:} \url{https://unicode.org}
\end{tecnico}

% =============================================================================
\section{IEEE 754 — Números con punto decimal}
% =============================================================================

Los enteros son fáciles de representar en binario, pero ¿cómo representar
$3.14$ o $0.001$? El estándar \textbf{IEEE 754} (1985) define el formato
que usan prácticamente todos los procesadores del mundo.

\subsection{La idea: notación científica en binario}

Igual que en notación científica decimal escribimos
$3.14 \times 10^0$ o $6.022 \times 10^{23}$,
en binario escribimos
\[
  \pm\; 1.\underbrace{b_1 b_2 \cdots b_{22}}_{\text{mantisa}} \;\times\; 2^{\,e}
\]
La parte entera siempre es 1 (en binario normalizado), así que no se
almacena — es un \textbf{bit implícito gratis}.

\subsection{Formato \texttt{float} (32 bits) y \texttt{double} (64 bits)}

\begin{center}
\begin{tabular}{lcccc}
  \toprule
  \textbf{Tipo} & \textbf{Bits totales} & \textbf{Signo} &
  \textbf{Exponente} & \textbf{Mantisa} \\
  \midrule
  \texttt{float}  & 32 & 1 & 8  & 23 \\
  \texttt{double} & 64 & 1 & 11 & 52 \\
  \bottomrule
\end{tabular}
\end{center}

\begin{itemize}
  \item \textbf{Signo}: 0 = positivo, 1 = negativo.
  \item \textbf{Exponente}: almacenado con un \textit{bias} (127 para
        \texttt{float}, 1023 para \texttt{double}) para evitar negativos.
  \item \textbf{Mantisa}: los bits después del punto del $1.xxx\ldots$
\end{itemize}

\subsection{¿Qué fracciones tienen representación exacta?}

En base 10, una fracción $\frac{1}{n}$ tiene representación decimal
\textbf{finita} si y solo si $n$ no tiene factores primos distintos de
\textbf{2} y \textbf{5} (los factores de 10).

\begin{center}
\begin{tabular}{lll}
  \toprule
  \textbf{Fracción} & \textbf{Decimal} & \textbf{¿Exacta?} \\
  \midrule
  $\frac{1}{2}$ & $0.5$               & Sí — $2$ divide a $10$ \\
  $\frac{1}{4}$ & $0.25$              & Sí — $4 = 2^2$ \\
  $\frac{1}{5}$ & $0.2$               & Sí — $5$ divide a $10$ \\
  $\frac{1}{3}$ & $0.\overline{3}$    & No — $3$ no divide a $10$ \\
  $\frac{1}{6}$ & $0.1\overline{6}$   & No — tiene factor $3$ \\
  $\frac{1}{7}$ & $0.\overline{142857}$ & No — $7$ no divide a $10$ \\
  \bottomrule
\end{tabular}
\end{center}

En base 2, el único factor primo es el \textbf{2}.
Por eso solo las fracciones de la forma $\frac{1}{2^k}$ son exactas;
todo lo demás es periódico.

\begin{center}
\begin{tabular}{lll}
  \toprule
  \textbf{Fracción} & \textbf{Binario} & \textbf{¿Exacta en base 2?} \\
  \midrule
  $\frac{1}{2}$ & $0.1_2$                   & Sí \\
  $\frac{1}{4}$ & $0.01_2$                  & Sí \\
  $\frac{1}{8}$ & $0.001_2$                 & Sí \\
  $\frac{1}{3}$ & $0.\overline{01}_2$        & No — patrón 01 se repite \\
  $\frac{1}{5}$ & $0.\overline{0011}_2$      & No — patrón 0011 se repite \\
  $\frac{1}{10}$& $0.0\overline{0011}_2$     & No ($\frac{1}{5}\div 2$) \\
  \bottomrule
\end{tabular}
\end{center}

\begin{recuerda}
  La representación periódica no es un defecto de la computadora:
  es una propiedad matemática. $\frac{1}{3}$ es igual de
  ``infinita'' en decimal como en binario. La diferencia es que en base 10
  lo vemos con $\frac{1}{3}$ pero no con $\frac{1}{10}$, y en base 2
  pasa lo contrario.
\end{recuerda}

\subsection{Cómo convertir la parte fraccionaria a binario}

El algoritmo es análogo a las divisiones sucesivas para la parte entera,
pero en sentido inverso: \textbf{multiplica por 2 repetidamente}.

\begin{enumerate}
  \item Multiplica la parte fraccionaria por 2.
  \item Si el resultado es $\geq 1$: escribe un \textbf{1} y conserva
        solo la parte fraccionaria (resta 1).
  \item Si el resultado es $< 1$: escribe un \textbf{0}.
  \item Repite con la parte fraccionaria restante.
        Si llegas a 0, terminaste. Si el estado se repite, el número
        es \textbf{periódico}.
\end{enumerate}

\begin{ejemplo}[Convertir $\frac{1}{4} = 0.25$ a binario]
\begin{center}
\begin{tabular}{r@{\ $\times 2=$\ }ll}
  $0.25$ & $0.50$ & $\to$ bit \textbf{0} \\
  $0.50$ & $1.00$ & $\to$ bit \textbf{1}, resto $0.00$ \\
\end{tabular}
\end{center}
$\frac{1}{4} = 0.01_2$ \checkmark
\end{ejemplo}

\begin{ejemplo}[Convertir $\frac{1}{3} \approx 0.333\ldots$ a binario]
\begin{center}
\begin{tabular}{r@{\ $\times 2=$\ }ll}
  $0.\overline{3}$ & $0.\overline{6}$ & $\to$ bit \textbf{0} \\
  $0.\overline{6}$ & $1.\overline{3}$ & $\to$ bit \textbf{1}, resto $0.\overline{3}$ \\
  $0.\overline{3}$ & $0.\overline{6}$ & $\to$ bit \textbf{0} \quad (¡mismo estado inicial!) \\
  $\vdots$ & $\vdots$ & $\vdots$ \\
\end{tabular}
\end{center}
El estado se repite desde el primer paso: $\frac{1}{3} = 0.\overline{01}_2$

\[
  \frac{1}{3} = 0.010101\ldots_2
\]
\end{ejemplo}

\begin{ejemplo}[Convertir $0.1_{10}$ a binario]
\begin{center}
\begin{tabular}{r@{\ $\times 2=$\ }ll}
  $0.1$ & $0.2$ & $\to$ bit \textbf{0} \\
  $0.2$ & $0.4$ & $\to$ bit \textbf{0} \\
  $0.4$ & $0.8$ & $\to$ bit \textbf{0} \\
  $0.8$ & $1.6$ & $\to$ bit \textbf{1}, resto $0.6$ \\
  $0.6$ & $1.2$ & $\to$ bit \textbf{1}, resto $0.2$ \\
  $0.2$ & $0.4$ & $\to$ bit \textbf{0} \quad (mismo estado que línea 2) \\
  $\vdots$ & $\vdots$ & $\vdots$ \\
\end{tabular}
\end{center}
$0.1_{10} = 0.0\overline{0011}_2$
\end{ejemplo}

\subsection{¿Por qué \texttt{0.1 + 0.2 $\neq$ 0.3} en la computadora?}

Porque $0.1_{10} = 0.0\overline{0011}_2$ es un número periódico en
binario. Al almacenarlo en un \texttt{float} (23 bits de mantisa) o
\texttt{double} (52 bits), la secuencia infinita se \textbf{trunca},
introduciendo un error pequeño pero no nulo.
Al sumar dos de esos errores el resultado no coincide exactamente con
la representación truncada de $0.3$.

\begin{advertencia}
  Nunca compares flotantes con \texttt{==} en C++. Usa una tolerancia:\\
  \texttt{|a - b| < 1e-9}. Ignorar esto es una fuente frecuente de
  bugs difíciles de rastrear.
\end{advertencia}

\begin{tecnico}[Valores especiales en IEEE 754]
  El estándar reserva patrones para casos extremos:
  \begin{itemize}
    \item \textbf{Infinito} ($\pm\infty$): exponente todo 1s, mantisa = 0.
          Resultado de dividir entre cero.
    \item \textbf{NaN} (\textit{Not a Number}): exponente todo 1s,
          mantisa $\neq$ 0. Resultado de $0/0$, $\sqrt{-1}$, etc.
    \item \textbf{Cero}: todo 0s (con signo, por eso existe $-0.0$).
  \end{itemize}
\end{tecnico}

% =============================================================================
\section{Ejercicios}
% =============================================================================

\subsection{Complemento a 2}

\begin{enumerate}

  \item Representa los siguientes valores en \textbf{8 bits en complemento
        a 2}. Muestra el proceso (invertir + sumar 1).
        \begin{enumerate}
          \item $-35$
          \item $-1$
          \item $-128$
        \end{enumerate}

  \item Dado el siguiente patrón de 8 bits interpretado como complemento
        a 2, ¿qué valor decimal representa? Muestra cómo lo determinas.
        \begin{enumerate}
          \item \texttt{1111 0110}
          \item \texttt{1000 0001}
          \item \texttt{0111 1111}
        \end{enumerate}

  \item ¿Cuál es el rango de valores representables con \textbf{16 bits}
        en complemento a 2? ¿Y sin signo?

  \item Un programador declara una variable de 8 bits sin signo
        (\texttt{uint8\_t}) y le asigna el valor 200. Luego la niega
        con el algoritmo de complemento a 2. ¿Qué valor obtiene?
        ¿Por qué el resultado es sorprendente?

\end{enumerate}

\subsection{ASCII}

\begin{enumerate}[resume]

  \item Decodifica el siguiente mensaje. Cada byte está en hexadecimal.
        \[
          \texttt{4C\ 55\ 4E\ 41\ 54\ 49\ 43}
        \]

  \item Escribe tu nombre completo (sin acentos) como secuencia de bytes
        en hexadecimal. Indica cuántos bytes ocupa.

  \item El byte \texttt{0x6C} es una letra minúscula. Sin consultar la
        tabla, determina cuál es la mayúscula correspondiente y su valor
        hexadecimal. Explica el procedimiento.

  \item ¿Qué imprimiría un programa que lee los bytes
        \texttt{0x39\ 0x0A\ 0x38\ 0x0A} y los envía a la pantalla
        carácter por carácter?

\end{enumerate}

\subsection{Fracciones en binario}

\begin{enumerate}[resume]

  \item Convierte las siguientes fracciones a binario usando el algoritmo
        de multiplicación por 2. Indica si el resultado es exacto o
        periódico y, en ese caso, cuál es el patrón que se repite.
        \begin{enumerate}
          \item $\dfrac{3}{8}$
          \item $\dfrac{1}{6}$
          \item $0.625$
        \end{enumerate}

  \item Sin hacer la conversión, determina cuáles de las siguientes
        fracciones tienen representación \textbf{exacta} en binario y
        cuáles no. Justifica con los factores primos.
        \[
          \frac{1}{16}, \quad \frac{1}{12}, \quad \frac{3}{4},
          \quad \frac{1}{14}, \quad \frac{5}{8}
        \]

  \item Convierte $5.375_{10}$ a binario (parte entera y fraccionaria).

\end{enumerate}

\subsection{IEEE 754}

\begin{enumerate}[resume]

  \item Clasifica cada número como exactamente representable o no en
        \texttt{float}. Justifica.
        \begin{enumerate}
          \item $0.5$
          \item $0.1$
          \item $\frac{1}{32}$
          \item $\frac{1}{10}$
        \end{enumerate}

  \item Un \texttt{float} de 32 bits tiene el siguiente patrón:
        \[
          \underbrace{0}_{signo}\;
          \underbrace{01111111}_{exponente}\;
          \underbrace{10000000000000000000000}_{mantisa}
        \]
        \begin{enumerate}
          \item ¿Es positivo o negativo?
          \item El exponente real es $e - 127$. ¿Cuánto vale aquí?
          \item La mantisa representa $1.1_2$. ¿Qué número decimal es?
        \end{enumerate}

  \item Explica con tus palabras por qué la expresión \texttt{0.1 + 0.2 == 0.3}
        devuelve \texttt{false} en C++, aunque matemáticamente sea correcta.

\end{enumerate}

% =============================================================================
\section{Resumen}
% =============================================================================

\begin{recuerda}
\begin{itemize}
  \item Con $n$ bits hay $2^n$ patrones; sin signo cubren 0 a $2^n-1$.
  \item \textbf{Complemento a 2}: rango $-2^{n-1}$ a $+2^{n-1}-1$.
        Para negar: invertir todos los bits y sumar 1.
  \item \texttt{11111111} (8 bits) $= -1$; \quad
        \texttt{10000000} (8 bits) $= -128$.
  \item \textbf{ASCII}: cada carácter es un byte.
        Mayúsculas = \texttt{0x41}–\texttt{0x5A};
        minúsculas = mayúscula $+ 32$.
  \item $\frac{1}{n}$ es exacta en base $b$ solo si $n$ no tiene factores
        primos ajenos a $b$. En base 2: solo $\frac{1}{2^k}$ es exacta.
  \item Para convertir una parte fraccionaria a binario: multiplica por 2,
        anota el dígito entero, repite con el resto.
  \item \textbf{IEEE 754}: flotantes = signo + exponente + mantisa.
        Los errores de redondeo son normales; nunca compares con \texttt{==}.
\end{itemize}
\end{recuerda}

\end{document}
